% =====================================================================
%  ENTORNOS PROPIOS
%  Los seis cuadros que estructuran cada capítulo. Úsalos siempre igual:
%  el lector debe poder hojear el libro buscando solo los recuadros.
% =====================================================================
\usepackage[most]{tcolorbox}
\tcbuselibrary{breakable,skins}

% --- Plantilla base ---------------------------------------------------
\tcbset{
  cursobase/.style={
    breakable, enhanced, sharp corners=downhill,
    boxrule=0pt, leftrule=3pt, toprule=0pt, bottomrule=0pt, rightrule=0pt,
    left=8pt, right=8pt, top=6pt, bottom=6pt,
    fonttitle=\sffamily\bfseries\small,
    coltitle=black, attach title to upper=\par\medskip
  }
}

% --- 1. Idea clave ----------------------------------------------------
\newtcolorbox{clave}[1][]{cursobase,
  colback=fusalClaro, colframe=fusalPrim,
  title={IDEA CLAVE}, #1}

% --- 2. Caso real del equipo -----------------------------------------
\newtcolorbox{caso}[1][]{cursobase,
  colback=fusalPrim!5, colframe=fusalPrim!70,
  title={CASO FUSAL\ifx\relax#1\relax\else\ --- #1\fi}}

% --- 3. Referencia al reglamento -------------------------------------
\newtcolorbox{regla}[1][]{cursobase,
  colback=gray!8, colframe=fusalGris,
  title={DEL REGLAMENTO\ifx\relax#1\relax\else\ --- #1\fi}}

% --- 4. Errores frecuentes -------------------------------------------
\newtcolorbox{ojo}[1][]{cursobase,
  colback=fusalSec!5, colframe=fusalSec,
  title={OJO --- ERRORES FRECUENTES}, #1}

% --- 5. Entregable de la sesión --------------------------------------
\newtcolorbox{entregable}[1][]{cursobase,
  colback=green!5, colframe=green!45!black,
  title={ENTREGABLE}, #1}

% --- 6. Preguntas de los jueces --------------------------------------
\newtcolorbox{juez}[1][]{cursobase,
  colback=orange!6, colframe=orange!70!black,
  title={LO QUE TE VA A PREGUNTAR UN JUEZ}, #1}

% --- Objetivos de aprendizaje (cabecera de cada capítulo) -------------
\newenvironment{objetivos}
  {\begin{tcolorbox}[cursobase, colback=white, colframe=fusalPrim,
      title={AL TERMINAR ESTE CAPÍTULO DEBES SABER}]
   \begin{itemize}[leftmargin=*,label=\textcolor{fusalPrim}{\textbullet}]}
  {\end{itemize}\end{tcolorbox}}

% --- Ficha de sesión (para impartir el capítulo como clase) -----------
% Uso: \fichasesion{Duración}{Requisitos previos}{Material necesario}
\newcommand{\fichasesion}[3]{%
  \begin{tcolorbox}[cursobase, colback=gray!5, colframe=fusalGris,
      title={FICHA DE SESIÓN}]
  \begin{description}[leftmargin=!,labelwidth=2.6cm,font=\sffamily\bfseries\small]
    \item[Duración] #1
    \item[Requisitos] #2
    \item[Material] #3
  \end{description}
  \end{tcolorbox}}
